% Options for packages loaded elsewhere
\PassOptionsToPackage{unicode}{hyperref}
\PassOptionsToPackage{hyphens}{url}
%
\documentclass[
]{article}
\usepackage{amsmath,amssymb}
\usepackage{iftex}
\ifPDFTeX
  \usepackage[T1]{fontenc}
  \usepackage[utf8]{inputenc}
  \usepackage{textcomp} % provide euro and other symbols
\else % if luatex or xetex
  \usepackage{unicode-math} % this also loads fontspec
  \defaultfontfeatures{Scale=MatchLowercase}
  \defaultfontfeatures[\rmfamily]{Ligatures=TeX,Scale=1}
\fi
\usepackage{lmodern}
\ifPDFTeX\else
  % xetex/luatex font selection
\fi
% Use upquote if available, for straight quotes in verbatim environments
\IfFileExists{upquote.sty}{\usepackage{upquote}}{}
\IfFileExists{microtype.sty}{% use microtype if available
  \usepackage[]{microtype}
  \UseMicrotypeSet[protrusion]{basicmath} % disable protrusion for tt fonts
}{}
\makeatletter
\@ifundefined{KOMAClassName}{% if non-KOMA class
  \IfFileExists{parskip.sty}{%
    \usepackage{parskip}
  }{% else
    \setlength{\parindent}{0pt}
    \setlength{\parskip}{6pt plus 2pt minus 1pt}}
}{% if KOMA class
  \KOMAoptions{parskip=half}}
\makeatother
\usepackage{xcolor}
\usepackage[margin=2.5cm]{geometry}
\usepackage{listings}
\newcommand{\passthrough}[1]{#1}
\lstset{defaultdialect=[5.3]Lua}
\lstset{defaultdialect=[x86masm]Assembler}
\usepackage{longtable,booktabs,array}
\usepackage{calc} % for calculating minipage widths
% Correct order of tables after \paragraph or \subparagraph
\usepackage{etoolbox}
\makeatletter
\patchcmd\longtable{\par}{\if@noskipsec\mbox{}\fi\par}{}{}
\makeatother
% Allow footnotes in longtable head/foot
\IfFileExists{footnotehyper.sty}{\usepackage{footnotehyper}}{\usepackage{footnote}}
\makesavenoteenv{longtable}
\setlength{\emergencystretch}{3em} % prevent overfull lines
\providecommand{\tightlist}{%
  \setlength{\itemsep}{0pt}\setlength{\parskip}{0pt}}
\setcounter{secnumdepth}{-\maxdimen} % remove section numbering
\ifLuaTeX
\usepackage[bidi=basic]{babel}
\else
\usepackage[bidi=default]{babel}
\fi
\babelprovide[main,import]{french}
% get rid of language-specific shorthands (see #6817):
\let\LanguageShortHands\languageshorthands
\def\languageshorthands#1{}
\ifLuaTeX
  \usepackage{selnolig}  % disable illegal ligatures
\fi
\usepackage{bookmark}
\IfFileExists{xurl.sty}{\usepackage{xurl}}{} % add URL line breaks if available
\urlstyle{same}
\hypersetup{
  pdflang={fr},
  hidelinks,
  pdfcreator={LaTeX via pandoc}}

\author{}
\date{}

\begin{document}

\section{TP --- Télécharger, stocker et charger un dataset dans Google
Colab}\label{tp-tuxe9luxe9charger-stocker-et-charger-un-dataset-dans-google-colab}

\textbf{Objectif :} télécharger un jeu de données depuis Kaggle, le
Fonds monétaire international (FMI) ou le Haut-Commissariat au Plan
(HCP) du Maroc, le stocker dans Google Drive ou sur GitHub, puis le
charger dans Google Colab avec Python et Pandas et l'exploiter avec
SQLite.

\subsection{1. Architecture générale}\label{architecture-guxe9nuxe9rale}

\begin{lstlisting}
Kaggle / FMI / HCP
       |
   CSV ou Excel
       |
  +----+----+
  |         |
Google Drive  GitHub (dépôt public)
  |         |
  +----+----+
       |
 Google Colab + Pandas
       |
 SQLite + requêtes SQL
\end{lstlisting}

\subsection{2. Télécharger un
dataset}\label{tuxe9luxe9charger-un-dataset}

\begin{itemize}
\tightlist
\item
  \textbf{Kaggle :} \url{https://www.kaggle.com/datasets} --- rechercher
  un dataset, ouvrir sa page et cliquer sur « Download ». Un compte peut
  être nécessaire.
\item
  \textbf{FMI :} \url{https://www.imf.org/en/Data} --- sélectionner une
  base et des indicateurs, puis exporter les données si l'option existe.
\item
  \textbf{HCP Maroc :} \url{https://www.hcp.ma} --- rechercher une
  publication statistique et télécharger un tableau CSV ou Excel
  lorsqu'il est disponible.
\end{itemize}

Dans la suite, on suppose que le fichier s'appelle
\passthrough{\lstinline!dataset.csv!}. Vérifier les conditions de
licence avant toute republication.

\subsection{3. Méthode A --- Stocker le dataset sur Google
Drive}\label{muxe9thode-a-stocker-le-dataset-sur-google-drive}

\subsubsection{Étape A1 --- Importer le
fichier}\label{uxe9tape-a1-importer-le-fichier}

\begin{enumerate}
\def\labelenumi{\arabic{enumi}.}
\tightlist
\item
  Ouvrir \url{https://drive.google.com}.
\item
  Créer un dossier \passthrough{\lstinline!TP\_SQL\_Dataset!}.
\item
  Importer \passthrough{\lstinline!dataset.csv!} dans ce dossier.
\item
  Ouvrir un notebook sur \url{https://colab.research.google.com}.
\end{enumerate}

\subsubsection{Étape A2 --- Monter Google
Drive}\label{uxe9tape-a2-monter-google-drive}

\begin{lstlisting}[language=Python]
from google.colab import drive

# Autoriser l'accès à Google Drive
drive.mount('/content/drive')
\end{lstlisting}

\subsubsection{Étape A3 --- Charger les
données}\label{uxe9tape-a3-charger-les-donnuxe9es}

\begin{lstlisting}[language=Python]
import pandas as pd
from pathlib import Path

chemin = Path('/content/drive/MyDrive/TP_SQL_Dataset/dataset.csv')

if not chemin.is_file():
    raise FileNotFoundError(f'Fichier introuvable : {chemin}')

df = pd.read_csv(chemin)
display(df.head(10))
print('Nombre de lignes :', df.shape[0])
print('Nombre de colonnes :', df.shape[1])
\end{lstlisting}

\textbf{Variantes :} pour Excel, utiliser
\passthrough{\lstinline!pd.read\_excel(chemin)!} avec un chemin
\passthrough{\lstinline!.xlsx!} ; pour un CSV séparé par des
points-virgules, utiliser
\passthrough{\lstinline!pd.read\_csv(chemin, sep=';')!}. Vérifier aussi
l'encodage du fichier.

\subsection{4. Méthode B --- Stocker le dataset sur
GitHub}\label{muxe9thode-b-stocker-le-dataset-sur-github}

\subsubsection{Étape B1 --- Créer un dépôt et téléverser le
fichier}\label{uxe9tape-b1-cruxe9er-un-duxe9puxf4t-et-tuxe9luxe9verser-le-fichier}

\begin{enumerate}
\def\labelenumi{\arabic{enumi}.}
\tightlist
\item
  Se connecter à \url{https://github.com} et choisir \textbf{New
  repository}.
\item
  Nommer le dépôt \passthrough{\lstinline!TP\_SQL\_Dataset!}.
\item
  Choisir \textbf{Public} si le chargement doit se faire sans
  authentification.
\item
  Créer le dépôt, puis choisir \textbf{Add file → Upload files}.
\item
  Ajouter \passthrough{\lstinline!dataset.csv!} et valider avec
  \textbf{Commit changes}.
\end{enumerate}

\textbf{Attention :} ne pas publier de données confidentielles ni de
dataset dont la licence interdit la redistribution. Les dépôts GitHub ne
sont pas adaptés aux très gros fichiers de données.

\subsubsection{Étape B2 --- Récupérer le lien
Raw}\label{uxe9tape-b2-ruxe9cupuxe9rer-le-lien-raw}

Ouvrir \passthrough{\lstinline!dataset.csv!} dans le dépôt, cliquer sur
\textbf{Raw} et copier son URL. Sa structure est généralement :

\begin{lstlisting}
https://raw.githubusercontent.com/UTILISATEUR/TP_SQL_Dataset/main/dataset.csv
\end{lstlisting}

Remplacer \passthrough{\lstinline!UTILISATEUR!} par le compte réel et
\passthrough{\lstinline!main!} par le nom de la branche utilisée si
nécessaire.

\subsubsection{Étape B3 --- Charger directement depuis
GitHub}\label{uxe9tape-b3-charger-directement-depuis-github}

\begin{lstlisting}[language=Python]
import pandas as pd

url = (
    'https://raw.githubusercontent.com/'
    'UTILISATEUR/TP_SQL_Dataset/main/dataset.csv'
)

df = pd.read_csv(url)
display(df.head(10))
print('Dimensions :', df.shape)
print('Colonnes :', df.columns.tolist())
\end{lstlisting}

Ce script est un \textbf{modèle à personnaliser} : l'URL d'exemple ne
pointe pas vers un dataset réel.

\subsection{5. Méthode C --- Script unique : choisir Drive ou
GitHub}\label{muxe9thode-c-script-unique-choisir-drive-ou-github}

Modifier \passthrough{\lstinline!SOURCE!} et les chemins en fonction du
fichier réellement utilisé.

\begin{lstlisting}[language=Python]
import pandas as pd

SOURCE = 'drive'  # 'drive' ou 'github'

if SOURCE == 'drive':
    from google.colab import drive
    drive.mount('/content/drive')
    fichier = '/content/drive/MyDrive/TP_SQL_Dataset/dataset.csv'
else:
    fichier = (
        'https://raw.githubusercontent.com/'
        'UTILISATEUR/TP_SQL_Dataset/main/dataset.csv'
    )

df = pd.read_csv(fichier)
display(df.head(10))
print('Dimensions :', df.shape)
\end{lstlisting}

\subsection{6. Exploiter ensuite le dataset avec
SQLite}\label{exploiter-ensuite-le-dataset-avec-sqlite}

Le DataFrame \passthrough{\lstinline!df!} a déjà été chargé par l'une
des méthodes précédentes.

\begin{lstlisting}[language=Python]
import sqlite3
import pandas as pd

# Base locale temporaire dans l'environnement Colab
conn = sqlite3.connect('tp_sql.db')

# Importer les données dans la table SQLite « donnees »
df.to_sql('donnees', conn, if_exists='replace', index=False)

# Première requête SQL
requete = '''
SELECT *
FROM donnees
LIMIT 10;
'''
resultat = pd.read_sql_query(requete, conn)
display(resultat)

# Vérifier le nombre de lignes importées
total = pd.read_sql_query(
    'SELECT COUNT(*) AS total FROM donnees;', conn
)
display(total)

conn.close()
\end{lstlisting}

\textbf{Persistance :} \passthrough{\lstinline!tp\_sql.db!} est créé ici
dans l'environnement temporaire de Colab. Pour le conserver, copier le
fichier vers Drive ou créer directement la base dans un dossier Drive
monté.

\subsection{7. Comparaison des deux
méthodes}\label{comparaison-des-deux-muxe9thodes}

\begin{longtable}[]{@{}
  >{\raggedright\arraybackslash}p{(\columnwidth - 4\tabcolsep) * \real{0.3333}}
  >{\raggedright\arraybackslash}p{(\columnwidth - 4\tabcolsep) * \real{0.3333}}
  >{\raggedright\arraybackslash}p{(\columnwidth - 4\tabcolsep) * \real{0.3333}}@{}}
\toprule\noalign{}
\begin{minipage}[b]{\linewidth}\raggedright
Critère
\end{minipage} & \begin{minipage}[b]{\linewidth}\raggedright
Google Drive
\end{minipage} & \begin{minipage}[b]{\linewidth}\raggedright
GitHub
\end{minipage} \\
\midrule\noalign{}
\endhead
\bottomrule\noalign{}
\endlastfoot
Facilité pour les débutants & Très bonne & Bonne \\
Stockage privé & Oui & Oui, avec authentification pour l'accès
programmatique \\
Chargement direct par URL & Pas avec le chemin Drive ordinaire & Oui,
avec Raw et dépôt public \\
Partage & Autorisations Drive & Dépôt et collaborateurs \\
Gestion des versions & Limitée & Très bonne \\
Données volumineuses & Généralement plus adapté & Limites de taille et
de dépôt \\
Usage pédagogique & TP individuels & TP collaboratifs et
reproductibilité \\
\end{longtable}

\subsection{8. Recommandation
pédagogique}\label{recommandation-puxe9dagogique}

Privilégier \textbf{Google Drive} comme méthode principale pour des
étudiants débutants. Présenter \textbf{GitHub} comme méthode
complémentaire pour aborder le partage des datasets, la reproductibilité
et le versionnement.

\textbf{Prolongement :} demander aux étudiants de construire 15 requêtes
SQL progressives sur le dataset chargé, puis de présenter et commenter
les résultats avec Pandas.

\end{document}
